\chapter{System Architecture}
\label{ch:architecture}

This chapter presents the architecture of \voname{} as a whole: the decomposition
into microservices, the clients that consume them, the gateway that guards them,
and the cross-cutting mechanisms --- authentication, synchronous and asynchronous
communication, real-time messaging, and persistence --- that make the parts work
as a coherent system. Chapter~\ref{ch:backend} then descends into each service in
detail; here the goal is the big picture and the patterns that repeat everywhere.

\section{Architectural Style}

\voname{} follows a \textbf{microservices} architecture with an \textbf{API
gateway} at the edge and an \textbf{event-driven} backbone for cross-service
communication. The guiding principles are:

\begin{description}
  \item[One capability per service.] Each service owns a single business
  capability --- identity, workspaces, chat, rooms, notifications --- and nothing
  else. A service is the sole owner of its data; no other service reaches into
  its database.
  \item[Security at the edge.] Authentication is verified once, at the gateway.
  Internal services trust identity that the gateway has already validated and
  forwarded, and never re-implement token parsing.
  \item[Synchronous where an answer is needed, asynchronous where it is not.]
  A service calls another synchronously only when it needs an immediate
  answer (for example, ``what is this user's role in this workspace?''). For
  changes that other services should merely \emph{react} to, it publishes an
  event and moves on.
  \item[Polyglot persistence.] Each service uses the database that best fits its
  data: a relational store for structured, transactional data; a document store
  for flexible, high-volume data; and an in-memory store for ephemeral state and
  caches.
\end{description}

\section{System Overview}

Figure~\ref{fig:sysarch} shows the full topology: three families of clients, the
gateway, the nine backend modules, the message broker, and the datastores.

\begin{figure}[htbp]
\centering
\begin{tikzpicture}[
  font=\footnotesize,
  client/.style={draw=voIndigo,fill=voIndigo!8,rounded corners=3pt,align=center,
                 minimum height=0.9cm,text width=2.5cm,inner sep=4pt},
  gw/.style={draw=voDark,fill=voDark!8,rounded corners=3pt,align=center,
             minimum height=0.9cm,text width=2.4cm},
  svc/.style={draw=voPrimary,fill=voPrimary!8,rounded corners=3pt,align=center,
              minimum height=0.85cm,text width=2.2cm,inner sep=3pt},
  scaf/.style={draw=voGray,fill=voLight,rounded corners=3pt,align=center,dashed,
              minimum height=0.85cm,text width=2.2cm,inner sep=3pt},
  infra/.style={draw=voAccent!70!black,fill=voAccent!10,rounded corners=3pt,
               align=center,minimum height=0.8cm,text width=2.0cm,inner sep=3pt},
  arr/.style={-{Stealth[]},voGray,thick},
  darr/.style={-{Stealth[]},voAccent!60!black,thick,dashed},
]
% clients
\node[client] (web) at (0,9)  {Web client\\\scriptsize React};
\node[client] (mob) at (3,9)  {Mobile client\\\scriptsize Flutter};
\node[client] (game) at (6.3,9) {2D Office\\\scriptsize Phaser + Colyseus};

% gateway
\node[gw] (gw) at (1.5,7.2) {API Gateway\\\scriptsize :8080};

% services row 1
\node[svc] (user) at (-1.2,5.4) {User\\\scriptsize :8091};
\node[svc] (wsx)  at (1.3,5.4)  {Workspace\\\scriptsize :8087};
\node[svc] (chat) at (3.8,5.4)  {Chat\\\scriptsize :8084};
\node[svc] (room) at (6.3,5.4)  {Room\\\scriptsize :8086};
% services row 2
\node[svc] (notif) at (-1.2,3.9) {Notifications\\\scriptsize :8082};
\node[scaf] (cal) at (1.3,3.9) {Calendar\\\scriptsize :8083 (stub)};
\node[scaf] (task) at (3.8,3.9) {Tasks\\\scriptsize :8085 (stub)};
\node[scaf] (shared) at (6.3,3.9) {Shared\\\scriptsize library};

% broker
\node[infra,fill=voIndigo!10,draw=voIndigo] (mq) at (9.4,4.65) {RabbitMQ\\\scriptsize event bus};

% datastores
\node[infra] (mysql) at (-1.2,2.0) {MySQL};
\node[infra] (pg) at (1.3,2.0) {PostgreSQL};
\node[infra] (mongo) at (3.8,2.0) {MongoDB};
\node[infra] (redis) at (6.3,2.0) {Redis};

% external
\node[client,draw=voAccent!70!black,fill=voAccent!8] (agora) at (9.4,9) {Agora\\\scriptsize media SFU};
\node[client,draw=voAccent!70!black,fill=voAccent!8] (smtp) at (9.4,7.4) {SMTP\\\scriptsize e-mail};

% client -> gateway
\draw[arr] (web) -- (gw); \draw[arr] (mob) -- (gw);
\draw[arr] (game) to[out=250,in=60] (room);
% gateway -> services
\draw[arr] (gw) -- (user); \draw[arr] (gw) -- (wsx);
\draw[arr] (gw) -- (chat); \draw[arr] (gw) -- (room);
% services -> datastores
\draw[arr] (user) -- (mysql); \draw[arr] (wsx) -- (pg);
\draw[arr] (chat) -- (mongo); \draw[arr] (room) -- (mongo);
\draw[arr] (notif) -- (mongo);
\draw[arr] (chat) -- (redis); \draw[arr] (room) -- (redis);
\draw[arr] (notif) -- (redis);
% events
\draw[darr] (user) to[out=0,in=180] (mq);
\draw[darr] (room) -- (mq); \draw[darr] (wsx) to[out=0,in=200] (mq);
\draw[darr] (mq.north) to[out=120,in=40] (chat.north east);
\draw[darr] (mq.south) to[out=-90,in=-12] (notif.south east);
% external
\draw[arr] (web) to[out=90,in=180] (agora);
\draw[arr] (room) -- (agora);
\draw[arr] (notif) -- (smtp);
% internal RestClient (role lookups) — arced above the service row, clear of the labels
\draw[arr,voPrimary,densely dotted] (chat.north) to[out=125,in=55]
      node[above,font=\tiny,voPrimary]{role} (wsx.north);
\draw[arr,voPrimary,densely dotted] (room.north) to[out=130,in=50] (wsx.north);
\end{tikzpicture}
\caption{\voname{} system architecture. Solid grey arrows: request/response and
storage; dashed green: asynchronous events over RabbitMQ; dotted blue:
synchronous internal role lookups; dashed boxes: unimplemented scaffolds.}
\label{fig:sysarch}
\end{figure}

\begin{notebox}
\textbf{Honest note on maturity.} Of the nine backend modules, five are
implemented services (user, workspace, chat, room, notifications), two are
unimplemented Maven scaffolds (\svc{calendar} and \svc{tasks} --- only a POM,
properties, and README each), and one (\svc{shared-library}) is a placeholder
module with no code. This is stated plainly wherever these modules appear, so
that the architecture is described as it is, not as intended.
\end{notebox}

\section{Service Catalogue}

Table~\ref{tab:services} catalogues the backend modules with their port,
responsibility, persistence, and messaging role. Two facts are worth flagging
because they correct common assumptions: the event bus is \textbf{RabbitMQ}
(there is no Kafka anywhere in the system), and internal service-to-service HTTP
uses Spring's \textbf{\code{RestClient}} (not Feign, despite an unused Feign
dependency lingering on one classpath).

\begin{table}[htbp]
\centering
\caption{Backend service catalogue.}
\label{tab:services}
\scriptsize
\begin{tabularx}{\textwidth}{L{2.2cm} C{0.9cm} X L{2.0cm} L{2.2cm}}
\toprule
\textbf{Service} & \textbf{Port} & \textbf{Responsibility} & \textbf{Datastore} & \textbf{Messaging} \\
\midrule
Gateway API & 8080 & Edge routing, JWT validation, header injection & --- & --- \\
User & 8091 & Registration, login, JWT issuance, profile, OTP & MySQL & Produces \code{Notification\-Event} \\
Workspace & 8087 & Workspaces, desks, teams, invitations, map layout, roles & PostgreSQL + Flyway & Produces \code{Workspace\-ChannelEvent} \\
Chat & 8084 & Channels, DMs, messages, threads, read receipts, real-time & MongoDB + Redis & Consumes room \& workspace events \\
Room & 8086 & Voice/video rooms, presence, Agora tokens, proximity voice & MongoDB + Redis & Produces \code{Room\-ChannelEvent} \\
Notifications & 8082 & E-mails + in-app notifications, real-time push & MongoDB + Redis & Consumes \code{Notification\-Event} \\
Calendar & 8083 & \emph{(scaffold)} scheduling / meetings & --- & --- \\
Tasks & 8085 & \emph{(scaffold)} task assignment / tracking & --- & --- \\
Shared library & --- & \emph{(placeholder)} intended common DTOs/utilities & --- & --- \\
\bottomrule
\end{tabularx}
\end{table}

\section{The API Gateway and Security Model}
\label{sec:gateway}

Every client request enters through the \textbf{gateway} (Spring Cloud Gateway,
reactive). The gateway is the system's single security trust boundary, and its
filter runs before routing with three responsibilities.

\begin{enumerate}
  \item \textbf{Anti-spoofing.} It strips any client-supplied \code{X-User-Id}
  and \code{X-User-Role} headers from every incoming request. This is what makes
  it \emph{safe} for internal services to trust those headers later --- they can
  only ever be set by the gateway.
  \item \textbf{Edge isolation of internal endpoints.} Any path under
  \route{/api/internal/**} is rejected with \code{403} at the edge, so
  server-to-server endpoints are unreachable from outside.
  \item \textbf{Authentication.} Public paths (\route{/api/auth/*},
  \route{/api/otp/*}, and any \route{*/health}) pass through untouched. Every
  other request must carry a valid \code{Bearer} JWT; the gateway verifies the
  HS256 signature and expiry, extracts the \code{userId} claim, and injects
  \code{X-User-Id} (and a default \code{X-User-Role}) before forwarding.
\end{enumerate}

Figure~\ref{fig:authflow} shows the end-to-end authentication flow, from
credential exchange to an authenticated call reaching a downstream service.

\begin{figure}[htbp]
\centering
\begin{tikzpicture}[
  font=\footnotesize,
  node distance=0.0cm,
  lifeline/.style={voGray,densely dashed},
  actorbox/.style={draw=voPrimary,fill=voPrimary!8,rounded corners=2pt,
                   minimum width=2.0cm,minimum height=0.7cm,align=center},
  msg/.style={-{Stealth[]},voDark,thick},
  ret/.style={-{Stealth[]},voGray,densely dashed},
  note/.style={font=\scriptsize\itshape,voGray},
]
% actors
\node[actorbox] (c) at (0,7) {Client};
\node[actorbox] (u) at (4,7) {User Svc\\:8091};
\node[actorbox] (g) at (8,7) {Gateway\\:8080};
\node[actorbox] (s) at (12,7) {Chat Svc\\:8084};
% lifelines
\foreach \n in {c,u,g,s}{\draw[lifeline] (\n.south) -- ++(0,-6.3);}
% messages
\draw[msg] (0,6.3) -- node[above,note]{login (email, password)} (4,6.3);
\draw[ret] (4,5.7) -- node[above,note]{JWT (HS256, userId, 24h)} (0,5.7);
\draw[msg] (0,5.0) -- node[above,note]{request + Bearer JWT} (8,5.0);
\node[note,align=left] at (8.9,4.4) {strip X-User-*,};
\node[note,align=left] at (8.7,4.1) {verify JWT};
\draw[msg] (8,3.6) -- node[above,note]{+ X-User-Id / X-User-Role} (12,3.6);
\node[note,align=left] at (12.9,3.0) {trust headers,};
\node[note,align=left] at (13.0,2.7) {no JWT parse};
\draw[ret] (12,2.2) -- node[above,note]{response} (0,2.2);
% internal call
\draw[msg,voPrimary] (12,1.5) to[out=200,in=340] node[below,note,voPrimary]{internal role lookup + X-Internal-Token} (4.2,1.5);
\end{tikzpicture}
\caption{End-to-end authentication and identity-propagation flow.}
\label{fig:authflow}
\end{figure}

Beyond the JWT, the system uses a second, coarser credential for
\textbf{server-to-server} calls: internal endpoints (workspace role lookups, the
Colyseus session/presence API) are authenticated with a shared
\code{X-Internal-Token}. Because such endpoints are blocked at the gateway, this
token only ever travels inside the trusted network. Together, the two mechanisms
give a clean separation: \emph{user} identity via JWT at the edge, and
\emph{service} identity via a shared token behind it.

\section{Communication Patterns}

\subsection{Synchronous: internal REST}

When a service needs an immediate answer from another, it makes a direct HTTP
call using Spring's \code{RestClient}. In \voname{} the dominant example is
authorisation: both the chat and room services must know a caller's role
\emph{within a specific workspace} before allowing an action, so they call the
workspace service's internal endpoint
\route{/api/internal/workspace/\{id\}/members/\{userId\}/role}. These calls are
designed to \textbf{fail closed}: if the workspace service is unreachable, the
caller denies the action (returning \code{503}) rather than allowing it, so a
dependency outage cannot become a security hole.

\subsection{Asynchronous: the RabbitMQ event bus}

Most cross-service communication is asynchronous and event-driven. A service that
changes state publishes an \emph{event} to a RabbitMQ \code{DirectExchange}; a
consuming service, bound to a durable queue by a routing key, reacts to it. The
producer neither knows nor waits for the consumer. Table~\ref{tab:events} lists
every event flow in the system.

\begin{table}[htbp]
\centering
\caption{Event flows over RabbitMQ. All use durable direct exchanges, JSON
payloads, and consumer-side retry with exponential backoff.}
\label{tab:events}
\scriptsize
\begin{tabularx}{\textwidth}{L{2.0cm} L{2.4cm} L{2.6cm} L{2.0cm} X}
\toprule
\textbf{Producer} & \textbf{Exchange} & \textbf{Routing key} & \textbf{Consumer} & \textbf{Event / effect} \\
\midrule
User & \code{notifications.\-exchange} & \code{routing.key} & Notifications & \code{NotificationEvent}: send e-mail / in-app notice \\
Room & \code{room.exchange} & \code{room.channel.event} & Chat & \code{RoomChannelEvent}: create/sync a \code{ROOM} channel \\
Workspace & \code{workspace.\-exchange} & \code{workspace.\-channel.event} & Chat & \code{WorkspaceChannelEvent}: maintain the canonical workspace channel \\
\bottomrule
\end{tabularx}
\end{table}

This design has three payoffs that recur through the book. First,
\textbf{decoupling}: the room service can create a room and return immediately,
while the chat service materialises the bound channel a moment later. Second,
\textbf{resilience}: if the chat service is momentarily down, its events wait
durably in the queue and are processed on recovery. Third, \textbf{idempotency}:
every event carries a unique \code{eventId}, and consumers check it before
acting, so a redelivered message is harmless. The price is \textbf{eventual
consistency} --- a freshly created room's channel may \code{404} for a brief
window until its event is processed --- which clients are designed to tolerate
(``setting up chat\ldots'').

\section{The Real-time Layer}
\label{sec:realtime-arch}

Three services --- chat, room, and notifications --- expose real-time
functionality over STOMP-on-WebSocket. They share two patterns so consistently
that it is worth describing them once, centrally.

\subsection{Ticket-based WebSocket authentication}

Because a browser cannot set an \code{Authorization} header on a WebSocket
handshake, each real-time service authenticates the socket with a single-use
\emph{ticket}, illustrated in Figure~\ref{fig:wsticket}:

\begin{enumerate}
  \item The client makes an authenticated REST call, \route{POST .../ws-ticket}.
  The service generates a random ticket and stores it in Redis, mapped to the
  caller's identity, with a 60-second time-to-live.
  \item The client opens the WebSocket with the ticket as a query parameter,
  e.g. \route{/api/chat/connect?ticket=<uuid>}.
  \item A handshake interceptor atomically reads-and-deletes the ticket from
  Redis (\code{GETDEL}), binding the resulting identity to the STOMP session. The
  ticket cannot be reused.
\end{enumerate}

\begin{figure}[htbp]
\centering
\begin{tikzpicture}[font=\footnotesize,
  box/.style={draw=voPrimary,fill=voPrimary!8,rounded corners=2pt,align=center,
              minimum height=0.7cm,text width=2.2cm},
  redis/.style={draw=voAccent!70!black,fill=voAccent!10,rounded corners=2pt,
              align=center,minimum height=0.7cm,text width=1.8cm},
  a/.style={-{Stealth[]},thick,voDark},
  r/.style={-{Stealth[]},densely dashed,voGray},
  note/.style={font=\scriptsize\itshape,voGray},
]
\node[box] (cl) at (0,0) {Client};
\node[box] (sv) at (5,0) {Real-time\\Service};
\node[redis] (rd) at (9.5,0) {Redis};
\draw[a] (cl.30) -- node[above,note]{1. POST ws-ticket (JWT)} (sv.150);
\draw[a] (sv.30) -- node[above,note]{store ticket, TTL 60s} (rd.150);
\draw[r] (sv.210) -- node[below,note]{2. ticket} (cl.-30);
\draw[a] (cl) to[out=-40,in=200] node[below,note]{3. WS connect ?ticket=} (sv.-140);
\draw[a] (sv.-30) to[out=-30,in=210] node[below,note]{GETDEL (single-use)} (rd.-150);
\end{tikzpicture}
\caption{Ticket-based WebSocket handshake, used identically by the chat, room,
and notification services.}
\label{fig:wsticket}
\end{figure}

\subsection{The \texorpdfstring{\code{\{action, payload\}}}{action/payload} envelope}

Every real-time message the server pushes uses the same envelope: a JSON object
with an \code{action} string and a \code{payload}. A chat client, for instance,
receives \code{\{action:"NEW\_MESSAGE", payload:\{...\}\}},
\code{\{action:"EDIT\_MESSAGE", ...\}}, or \code{\{action:"TYPING", ...\}}; a room
client receives \code{PARTICIPANT\_JOINED}, \code{STATE\_CHANGED}, or
\code{ROOM\_CLOSED}; a notification client receives \code{NEW\_NOTIFICATION}. This
uniformity means every client can use one dispatch loop --- switch on
\code{action}, handle the \code{payload} --- regardless of service, which
noticeably simplifies the web and mobile code (Chapters~\ref{ch:web}
and~\ref{ch:mobile}).

STOMP \emph{destinations} follow an equally regular scheme: public broadcasts go
to \route{/topic/...} (a channel, thread, or room), a user's private stream goes
to \route{/user/queue/...} (errors, notifications), and clients send to
\route{/app/...} application destinations. Subscriptions to sensitive topics are
authorised by an interceptor that checks membership before allowing the
\code{SUBSCRIBE}.

\section{Polyglot Persistence}

Each service stores its data in the technology that fits it best
(Table~\ref{tab:services}):

\begin{description}
  \item[MySQL --- User service.] Accounts and verification requests are
  structured, relational, and transactional; a relational database with JPA is
  the natural home.
  \item[PostgreSQL with Flyway --- Workspace service.] The workspace domain is
  the most structurally rich in the system (workspaces, desks, teams, zones, map
  layers, tilesets, invitations) with many foreign-key relationships and
  optimistic-locking needs; PostgreSQL provides that, and Flyway migrations give
  the schema first-class version control.
  \item[MongoDB --- Chat, Room, Notifications.] Messages, rooms, and
  notifications are high-volume, append-heavy documents whose shape varies
  (mentions, attachments, arbitrary notification \code{data}); a document store
  fits, and each service uses its own database on the shared instance.
  \item[Redis --- Chat, Room, Notifications.] Ephemeral and latency-critical
  state lives in Redis: WebSocket tickets, read-receipt cursors, room presence
  and heartbeats, rate-limit counters, and idempotency claims. Keys are
  namespaced per service (\code{read:*}, \code{room:*}, \code{notif:*}) to share
  one instance safely.
\end{description}

\section{Technology Stack Summary}

Table~\ref{tab:stack} consolidates the technologies used across the platform, by
tier. The clients are covered in Chapters~\ref{ch:skyoffice}--\ref{ch:mobile}.

\begin{table}[htbp]
\centering
\caption{Technology stack by tier.}
\label{tab:stack}
\small
\begin{tabularx}{\textwidth}{L{3.2cm} X}
\toprule
\textbf{Tier} & \textbf{Technologies} \\
\midrule
Backend services & Java 21, Spring Boot (3.4.1--4.0.x), Spring Cloud Gateway, Spring Web/WebFlux, Spring Data JPA/Mongo/Redis, Spring AMQP, Spring WebSocket + STOMP, jjwt, MapStruct, Flyway \\
Data \& messaging & MySQL, PostgreSQL, MongoDB, Redis, RabbitMQ \\
Web client & React 18, MUI 7, Tailwind CSS, \code{@stomp/stompjs}, \code{agora-rtc-sdk-ng}, axios, React Router 7 \\
Mobile client & Flutter / Dart, Provider, \code{http}, \code{stomp\_dart\_client}, \code{awesome\_notifications}, \code{shared\_preferences} \\
2D office & TypeScript, Phaser 3, Colyseus, PeerJS/WebRTC, React + Redux Toolkit, Vite \\
External services & Agora (media SFU), Gmail SMTP (e-mail) \\
Build \& ops & Maven (multi-module), Docker Compose, health endpoints \\
\bottomrule
\end{tabularx}
\end{table}

\section{Deployment View}

For development, each service runs on its documented port on \code{localhost},
and its infrastructure dependencies (MySQL, PostgreSQL, MongoDB, Redis, RabbitMQ)
are brought up through Docker Compose definitions shipped with the services.
Spring Boot's Docker Compose integration lets a service start its own
dependencies automatically. In a containerised deployment, the gateway is the
only publicly exposed service; all others sit on an internal network, reachable
only through the gateway or, for internal endpoints, only by peers holding the
shared internal token. Service URLs are externalised as environment variables
(for example \code{CHAT\_SERVICE\_URL}), so the same images run unchanged from a
laptop to a cluster.

The chapters that follow now open each of these boxes in turn: the backend
services (Chapter~\ref{ch:backend}), the 2D office (Chapter~\ref{ch:skyoffice}),
and the web and mobile clients (Chapters~\ref{ch:web} and~\ref{ch:mobile}).
